% !Mode:: "TeX:UTF-8"
% VL 12, gebaut aus part-lernen/10-autograd-berechnungsgraph.md und
% part-architekturen-anwendungen/21-llms-benutzen.md
% (Dialogmuster und Zeiten: transcripts/vl11-wiederholung-llms-lokal-2025/notiz.md).
% !Mode:: "TeX:UTF-8"
\newcommand{\thepage}{} % to silence a hyperref warning about pdfpagelabels
\RequirePackage[]{silence}
\WarningsOff[hyperref]
\PassOptionsToPackage{log-declarations=false}{xparse}
%\RequirePackage[l2tabu, orthodox]{nag}
\RequirePackage{luatex85} % to make \pdf...= commands work with LuaTeX
\documentclass[aspectratio=169,t]{beamer}
\usepackage{ifthen}
\usepackage[T1]{fontenc}
\usepackage{iftex}
\ifLuaTeX
 \usepackage{fontspec}
 \setmainfont[Ligatures=TeX]{Source Sans Pro}
\else
 \usepackage[utf8]{inputenc}
\fi
\usepackage{german}
\usepackage{newunicodechar}
\newunicodechar{€}{\texteuro}
\newunicodechar{ß}{\ss}
\renewcommand{\SS}{\iffontchar\font"1E9E \symbol{"1E9E}\else SS\fi}
\usepackage{ellipsis}
\usepackage{microtype}
\usepackage{listings}
\usepackage{sourcecodepro}
\lstset{language=Python, basicstyle=\ttfamily}
\usepackage{amsfonts}
\usepackage{amssymb}
\usepackage{amsthm}
\usepackage{fancyhdr}
\usepackage{graphicx}

\newtheorem{proposition}[theorem]{Proposition}

\DeclareMathOperator*{\argmax}{argmax}
\DeclareMathOperator{\supp}{supp}
\newcommand{\E}{\ensuremath{\mathbb{E}}}
\newcommand{\V}{\ensuremath{\mathbb{V}}}

\setbeamercolor{emph}{fg=purple}
\renewcommand<>{\emph}[1]{%
  {\usebeamercolor[fg]{emph}\only#2{\itshape}#1}%
}



\hypersetup{draft=false,colorlinks=false,pdfpagemode=UseOutlines,plainpages=false,
            breaklinks=true,bookmarksnumbered=true,pdfborder=0 0 0,
            pdftitle={Einführung in Deep Learning WiSe 2026/27 HHU},
            pdfauthor={Konrad Voelkel},
            pdfsubject={data science}}

%\documentclass[aspectratio=169,t]{beamer}
\usetheme{hhubeamer} % put the hhu-beamer-theme in /usr/share/texmf/tex/latex/beamer/themes/

% -------------------- Code listings config --------------------
% Define custom colors for syntax highlighting
\definecolor{codegreen}{rgb}{0,0.6,0}
\definecolor{codegray}{rgb}{0.5,0.5,0.5}
\definecolor{codepurple}{rgb}{0.58,0,0.82}
\definecolor{backcolour}{rgb}{0.96,0.96,0.94} % Very light gray/beige
% Configure the listings package
\lstset{
    language=Python,
    backgroundcolor=\color{backcolour},   
    commentstyle=\color{codegreen}\itshape,
    keywordstyle=\color{blue}\bfseries,
    numberstyle=\tiny\color{codegray},
    stringstyle=\color{codepurple},
    basicstyle=\ttfamily\footnotesize, % Use a monospaced font
    breakatwhitespace=false,         
    breaklines=true,      
             % Enable automatic line breaking
    captionpos=b,                      % Caption at the bottom
    keepspaces=true,                 
    numbers=left,                      % Line numbers on the left
    numbersep=10pt,                    % Distance of line numbers from the code
    showspaces=false,                
    showstringspaces=false,            % Don't show weird underscores in strings
    showtabs=false,                  
    tabsize=4,
    frame=single,                      % Add a subtle border around the code
    rulecolor=\color{black!20},        % Border color
  }

\event[]{}
\author[Wintersemester 2026/27]{Dr. Konrad Völkel, Wintersemester 2026/27}

%\title{Data Science, VL 1}
%\date{16.04.2026}


%%% Local Variables:
%%% mode: latex
%%% TeX-master: t
%%% End:


% Fassung mit Dozentennotizen (zweiter Schirm): die nächste Zeile einkommentieren.
% \setbeameroption{show notes on second screen=right}

\graphicspath{{../images/}}

\title[Autograd und LLMs (VL 12)]{Autograd und LLMs benutzen\\\large Einführung in Deep Learning, VL 12}
\date{13.01.2027}

\begin{document}

\begin{frame}
  \titlepage

  \note{\textbf{Ablauf, 90 Minuten} (2025: 46 Min. Wiederholung, 50 Min. LLM-Praxis ---
    und 7 Min. Überzug, den wir uns dieses Jahr sparen):\\[2pt]
    \begin{tabular}{@{}rll@{}}
      0--4   & Folie 2       & Rückblick, heute zwei Hälften\\
      4--20  & Folien 3--6   & \textbf{Autograd-Notebook}: Graph, backward, step\\
      20--28 & Folien 7--8   & Freezing, geteilte Gewichte\\
      28--40 & Folien 9--10  & Zehn Wege, ein LLM zu nutzen\\
      40--48 & Folien 11--12 & Wie viel VRAM? Llamas\\
      48--60 & Folie 13      & \textbf{Ollama live im Terminal}\\
      60--74 & Folien 14--15 & Die Verteilung direkt auslesen (llama.cpp)\\
      74--84 & Folie 16      & Constrained Generation, JSON\\
      84--90 & Folie 17      & Übungsblatt\\
    \end{tabular}\\[4pt]
    Zwei Notebooks: \texttt{VL\_11\_backprop.ipynb} und \texttt{VL\_11\_llamacpp.ipynb}.
    Beide zeigen, nicht vorlesen. 2025 hat der Dozent hier überzogen --- notfalls die
    Wege-Liste kürzen.}
\end{frame}

\begin{frame}
  \frametitle{Heute zwei Hälften}

  \begin{center}
    \Large
    Erst: Was macht \texttt{loss.backward()} wirklich?\\[16pt]
    Dann: ein LLM ohne API benutzen.
  \end{center}

  \note{Die erste Hälfte ist Wiederholung von VL 6 --- diesmal aber im Code statt auf
    Papier. Das ist die Gelegenheit, Backpropagation ein zweites Mal zu sehen, wenn es
    beim ersten Mal zu schnell ging.}
\end{frame}

\begin{frame}[fragile]
  \frametitle{Ein sehr kleines Netz}

  \begin{lstlisting}
import torch

x = torch.tensor(3.0)                      # Input (Konstante)
w = torch.tensor(2.0, requires_grad=True)  # "Beobachte mich!"
b = torch.tensor(1.0, requires_grad=True)
z = w * x     # 6
y = z + b     # 7
L = y ** 2    # 49
  \end{lstlisting}

  \note{\texttt{requires\_grad} ist die ganze Magie: Damit merkt sich PyTorch, dass dieser
    Tensor beobachtet wird.
    Der Berechnungsgraph (ein DAG) entsteht \emph{beim Ausführen}, nicht beim Schreiben ---
    das ist der Unterschied zwischen Formel und Rechenreihenfolge, den wir in VL 6 hatten.}
\end{frame}

\begin{frame}[fragile]
  \frametitle{Der Graph ist sichtbar}

  \begin{lstlisting}
print(L.grad_fn)                       # <PowBackward0>
print(L.grad_fn.next_functions[0][0])  # <AddBackward0>
  \end{lstlisting}

  \bigskip
  \begin{center}
    \normalsize Jede \texttt{grad\_fn} weiß, wie man ihre Operation ableitet ---\\
    und hält die \emph{saved tensors}.
  \end{center}

  \note{Das sind genau die $z_k$ aus VL 6, die wir im Forward Pass gespeichert haben.
    \texttt{AccumulateGrad} sind die Blätter des Graphen, also die Parameter.
    Im Notebook zeigen und einmal den Graphen entlanglaufen --- das macht die Sache
    greifbar.}
\end{frame}

\begin{frame}[fragile]
  \frametitle{backward()}

  \begin{center}
    \Large
    $\dfrac{dL}{dw} = 2y \cdot 1 \cdot x = 14 \cdot 3 = 42$
  \end{center}

  \begin{lstlisting}
L.backward()
print(w.grad, b.grad)   # tensor(42.), tensor(14.)
  \end{lstlisting}

  \note{\texttt{backward()} ist der rekursive Aufruf der Kettenregel, initial mit
    $dL/dL = 1$. Die Zahlen stimmen --- nachgerechnet.
    Das ist derselbe Rückwärtsdurchlauf wie in VL 6, nur dass ihn niemand von Hand
    aufschreiben muss.}
\end{frame}

\begin{frame}[fragile]
  \frametitle{Und der Optimizer-Schritt}

  \begin{lstlisting}
with torch.no_grad():     # sonst wird der Update Teil des Graphen
    w -= lr * w.grad      # 2.0 - 0.01 * 42 = 1.58
    b -= lr * b.grad      # 1.0 - 0.01 * 14 = 0.86
    w.grad.zero_()        # sonst akkumuliert PyTorch
    b.grad.zero_()
  \end{lstlisting}

  \bigskip
  \begin{center}
    \normalsize Mehr macht \texttt{optimizer.step()} nicht.
  \end{center}

  \note{Zwei Fallen auf einer Folie: \texttt{no\_grad}, sonst landet das Update selbst im
    Graphen --- und \texttt{zero\_()}, sonst summiert PyTorch die Gradienten auf.
    Das erklärt die Zeile \texttt{optimizer.zero\_grad()} aus VL 4, über die sich alle
    wundern.
    Übrigens ist das Akkumulieren kein Bug: Man nutzt es für große Batches, die nicht in
    den Speicher passen.}
\end{frame}

\begin{frame}
  \frametitle{Freezing}

  \begin{center}
    \Large \texttt{w.requires\_grad = False}
  \end{center}

  \bigskip
  \begin{center}
    \normalsize Der Graph ist an dieser Stelle abgeschnitten --- der Optimizer
    ändert $w$ nicht mehr.
  \end{center}

  \note{So friert man beim Transfer Learning die Schichten eines vortrainierten Modells
    ein. Spart nebenbei Speicher, weil für eingefrorene Parameter keine Gradienten
    gehalten werden.
    Das ist der Mechanismus hinter „nur den letzten Layer nachtrainieren“.}
\end{frame}

\begin{frame}[fragile]
  \frametitle{Geteilte Gewichte}

  \begin{lstlisting}
x = torch.tensor([1.0, 2.0, 3.0])
k = torch.tensor([0.5], requires_grad=True)
o1, o2, o3 = x[0] * k, x[1] * k, x[2] * k   # wie nn.Conv1d
loss = o1 + o2 + o3
loss.backward()
print(k.grad)     # tensor([6.])
  \end{lstlisting}

  \bigskip
  \begin{center}
    \normalsize $1 + 2 + 3 = 6$: alle Verwendungen werden \emph{summiert}.
  \end{center}

  \note{Das ist die „Summe über alle Pfade“ aus VL 5, jetzt im Code: Im Graphen führen
    mehrere Kanten zu demselben Blatt, und die Beiträge addieren sich.
    Genau so funktioniert Backprop durch eine Faltungsschicht --- ein Kernel wird an jeder
    Position benutzt.}
\end{frame}

% ------------------------------------------------------------------ LLMs
\begin{frame}
  \frametitle{Zehn Wege, ein LLM zu nutzen}

  \begin{columns}[t]
    \begin{column}{0.48\textwidth}
      \begin{enumerate}\setlength{\itemsep}{1pt}
      \item Wortvektoren
      \item Satzvektoren
      \item die Verteilung auslesen
      \item Prompting
      \item selbst pretrainieren
      \end{enumerate}
    \end{column}
    \begin{column}{0.48\textwidth}
      \begin{enumerate}\setlength{\itemsep}{1pt}\setcounter{enumi}{5}
      \item Finetuning
      \item RLHF
      \item API
      \item Destillation
      \item Constrained Generation
      \end{enumerate}
    \end{column}
  \end{columns}

  \note{Am Beispiel-Task durchgehen: „Ist dieser Text positiv oder negativ gefärbt?“
    Weg 1 ist sehr billig, aber blind für längere Zusammenhänge; Weg 2 hat hohe Qualität,
    braucht aber ein LLM zur Inferenz (das ist das Übungsblatt dieser Woche); Weg 5 kostet für Llama 2 zwischen
    2 und 10 Mio. Dollar.
    Heute machen wir 1, 3 und 10 im Code.}
\end{frame}

\begin{frame}
  \frametitle{Wie viel Rechner braucht man?}

  \begin{center}
    \Large
    Llama 2 7B: $7 \cdot 4$ GB $= 28$ GB (32 Bit)\\[10pt]
    $7 \cdot 2$ GB $= 14$ GB (16 Bit)\\[10pt]
    quantisiert (4 Bit): unter 5 GB
  \end{center}

  \note{Die Faustformel ist die wichtigste Zahl der Sitzung: Parameterzahl mal Bytes pro
    Parameter.
    Llama2-7B ist etwa 8,6 GB groß und passt auf eine Consumer-Grafikkarte.
    Llama3.3-70B spielt in der Liga von GPT-4, braucht aber viel VRAM; Llama3.1-405B
    verteilt man auf mehrere GPUs.}
\end{frame}

\begin{frame}[fragile]
  \frametitle{Ollama}

  \begin{lstlisting}
curl -fsSL https://ollama.com/install.sh | sh
ollama run llama3.2:1b "Explain how LLMs work"
  \end{lstlisting}

  \bigskip
  \begin{center}
    \normalsize Baut auf llama.cpp / \texttt{ggml} auf.
  \end{center}

  \note{\textbf{Live im Terminal zeigen} --- das dauert zwei Minuten und beeindruckt mehr
    als jede Folie. Ein 1B-Modell läuft auf jedem Laptop.
    Danach im Browser das Chat-Interface zeigen.
    Wer es nachbauen will: Für die Übung reicht ein kleines Modell wie
    \texttt{qwen2.5:0.5b}, das sind 400 MB.}
\end{frame}

\begin{frame}[fragile]
  \frametitle{Weg 1: Embeddings}

  \begin{lstlisting}
import ollama, numpy as np

response = ollama.embeddings(model="qwen2.5:0.5b",
                             prompt="Heute ist es kalt draussen.")
vector = np.array(response["embedding"])
print(vector.shape, np.linalg.norm(vector))    # (896,)  ~300
  \end{lstlisting}

  \note{Nachgerechnet: Shape $(896,)$, Norm etwa 297.
    Das ist word2vec aus VL 9, nur mit einem viel größeren Modell --- und mit Kontext,
    weil der ganze Satz eingeht.
    Damit kann man semantisch suchen und clustern, ohne je Text zu generieren.}
\end{frame}

\begin{frame}[fragile]
  \frametitle{Weg 3: die Verteilung direkt}

  \lstset{basicstyle=\ttfamily\scriptsize}
  \begin{lstlisting}
probs = next_token_probs("The color of the sky is")
print(probs.shape)          # (151936,) -- die Vokabulargroesse
# ' changing' 10.3%, ' always' 6.7%, ' blue' 4.0%, ' a' 3.2%
print(f"{probs[red]:.2%}")  # 0.72%
  \end{lstlisting}

  \bigskip
  \begin{center}
    \normalsize Ein LLM liefert keinen Text, sondern eine \emph{Verteilung}.
  \end{center}

  \note{Das ist die Take-Home-Message der zweiten Hälfte. Der Text entsteht erst durch
    Sampling aus dieser Verteilung --- und was man statt \texttt{argmax} nimmt
    (Temperatur, Top-$k$, Top-$p$), ist eine eigene Entscheidung.
    Die Zahlen sind mit Qwen2.5-0.5B-Instruct Q4\_K\_M nachgerechnet; mit anderem Modell
    oder anderer Quantisierung sehen sie anders aus.
    Im Notebook \texttt{VL\_11\_llamacpp.ipynb} live zeigen.}
\end{frame}

\begin{frame}[fragile]
  \frametitle{Weg 10: Constrained Generation}

  \lstset{basicstyle=\ttfamily\scriptsize}
  \begin{lstlisting}
class Person(BaseModel):
    name: str; age: int; occupation: str; location: str

response = ollama.chat(model="qwen2.5:0.5b", format="json", messages=[...])
person = Person.model_validate_json(response["message"]["content"])
print(person)   # name='Sabine' age=29 occupation='doctor' location='Munich'
  \end{lstlisting}

  \note{Auch das ist nachgerechnet und kommt genau so heraus.
    Der Punkt: Statt in Prompts um ein Format zu \emph{betteln}, zwingt man das Sampling
    in ein Schema --- beim Generieren werden nur Token zugelassen, die das JSON gültig
    halten.
    Das ist der Baustein, auf dem agentische Workflows aufbauen.}
\end{frame}

\begin{frame}
  \frametitle{Übungsblatt}

  \begin{center}
    \Large BERT-Embeddings als Features für einen Klassifikator.
  \end{center}

  \bigskip
  \begin{center}
    Nächste Woche: GANs.
  \end{center}

  \note{Das ist Weg 2 aus der Liste. Das Notebook dazu liegt auf der Seite des Übungsblatts;
    es lädt ein bert-medium herunter und braucht \texttt{ipywidgets} für die
    Saliency-Zelle.
    Ansagen: Die Probeklausur kommt nächste Woche raus.}
\end{frame}

\end{document}

%%% Local Variables:
%%% mode: latex
%%% TeX-master: t
%%% End:
